\section{Limitations}\label{sec:limits}

\subsection{Threats to internal validity}

The robust hologram and the baseline are computed by the same engine, from the same random start of each seed, and scored by the same function on the same draws, so the comparison does not depend on a difference of implementation. The ablation uses a replica of the sigma-point loop that takes the design samples as an argument; the replica reproduces the module's hologram on the default configuration to \replicaDiff{} rad, and the evaluator of this paper reproduces the module's ensemble statistics to \evalCheckMax. The seeds are few: three seeds feed every part but the main comparison, so
the ranges of the ablation are ranges over three holograms and not confidence intervals, and a difference between two variants of less than the range is not a result. The bootstrap interval of a single seed measures the finite ensemble, and it is wide (\ciBelowN{} of \nSeedsAll{} intervals reach below the criterion). The ensemble seed of the acceptance test is one draw of 200; the five other ensembles of the held-out test give higher ratios, which suggests that the acceptance ensemble is not favourable but does not bound the variation of
the ensemble. The conjugate-gradient iterations end at a budget and not at a convergence criterion, and the optimizer is not a global one, so the holograms are those of one descent from one start. The first-order model that checks the theory uses finite differences of $10^{-3}$ waves, which is small against the phase offsets of the ensemble but not an analytical derivative.

\subsection{Threats to external validity}

The result is a simulation in a scalar Fourier model with an ideal device: no pixel crosstalk, no phase quantization, no fill factor, no flicker, no camera. The ablation of the design shows that the gain is tied to the match between the design prior and the ensemble in amplitude and in modes, and the ensemble here is drawn from the prior that the design assumes. A real residual is neither independent across modes nor Gaussian nor of one scale, varies across the field, and is a quantity that has to be measured; with the evaluation amplitude $1.5$ times the design amplitude the measured ratio is \misDMed{}, and with the design amplitude $1.5$ times the evaluation amplitude it is \ablPriorDMed. The array is 100 traps on 256 by 256 pixels, a small array compared with the arrays that motivate the work (about a thousand tweezers for the 256-atom machine~\cite{ebadi2021} and about twelve thousand sites for the 6,100-atom array~\cite{manetsch2025}); the cost of a run grows with the pixel count and the trap count (Section~\ref{sec:proposed-method}), and the benefit at that scale is not measured. The traps have equal weights, and the first-order
theory is stated for that case. The aberration enters as a phase offset on the SLM plane, as in the sources, and a field-dependent aberration, for which each trap sees a different wavefront, is not modelled.

\subsection{Threats to construct validity}

The spread $\sigma$ is the relative standard deviation of the trap intensities across one realization, and the 95th percentile is taken over realizations. The quantity that an experiment needs may be the trap depth of each site, its stability over time, or the loading probability, which depend on the intensity through physics outside this model; the spread of the intensities is a proxy. The efficiency counts the light in the traps and not the light in the zero order or in ghost spots. The ratio $\rho$ compares two 95th percentiles of 200 draws each, and
the 95th percentile of such a sample is a noisy statistic, which the bootstrap interval measures for one seed. The 100 traps of equal weight also make the spread a single number per realization; with unequal weights, or traps of different kinds, the construct would have to be defined per class of traps. The claim of the hypotheses H1 and H2 uses the thresholds of the product's acceptance test, and the thresholds are conventions, not requirements of an experiment.

\subsection{Scope}

The method applies to trap targets of point traps in two dimensions; a shaped field is refused. The sigma-point design covers aberrations; the displacement prior needs the sampled design, which is a different cost with different results, reported in Section~\ref{sec:results}. The method is an offline design: it needs on the order of \runSeconds{} seconds of wall-clock time per hologram in the conditions of the measurement and is not part of a live rearrangement loop. It is not a replacement for aberration correction, camera feedback or a model of the device, and it was not combined with any of them in these experiments. Combined use is plausible, since the correction fixes the nominal system around which the prior is centred and the device model fixes the forward map, but the interaction (a correction with its own error, a prior centred on an estimate, a cost evaluated on a device model with crosstalk) is not measured here and may change the optimum of the level and of the design amplitude.
